% reference_theme.tex — compact, monochrome serif reference presentation.
% Requires preamble.tex. Optional semantic blocks and page shells are separate.
% No class, page geometry, base font size, line spread, column count, direction,
% running furniture, title content or build policy is owned by this theme.
% <hllpdf:architecture-guard>
\ifdefined\hllpdfcoreloaded\else
  \errmessage{hllpdf: reference_theme.tex requires preamble.tex first}
  \expandafter\endinput
\fi
\ifdefined\hllpdfreaderthemeloaded
  \errmessage{hllpdf: reference_theme.tex conflicts with reader_theme.tex}
  \expandafter\endinput
\fi
\ifdefined\hllpdfacademicthemeloaded
  \errmessage{hllpdf: reference_theme.tex conflicts with academic_theme.tex}
  \expandafter\endinput
\fi
\ifdefined\hllpdfcleanthemeloaded
  \errmessage{hllpdf: reference_theme.tex conflicts with clean_theme.tex}
  \expandafter\endinput
\fi
\ifdefined\hllpdfcleanvariantloaded
  \errmessage{hllpdf: reference_theme.tex conflicts with clean_variant.tex}
  \expandafter\endinput
\fi
\ifdefined\hllpdfreferencethemeloaded
  \errmessage{hllpdf: reference_theme.tex loaded more than once}
  \expandafter\endinput
\fi
\def\hllpdfreferencethemeloaded{1}
% </hllpdf:architecture-guard>

% The core's pinned Hebrew Frank Ruhl Libre serif is retained. English uses
% the already verified TeX-distribution Latin Modern Roman faces/fallbacks.
\babelfont[english]{rm}[Ligatures=TeX,RawFeature={fallback=hebfb},
  BoldFeatures={RawFeature={fallback=hebfbbf}},
  BoldItalicFeatures={RawFeature={fallback=hebfbbf}}]{Latin Modern Roman}

% Professional selective numbering: only referenced equations receive an
% automatic number. Explicitly mark a central unreferenced equation through
% \noeqref{label}; use starred environments for intermediate calculations.
\mathtoolsset{showonlyrefs=true,showmanualtags=false}
\lakitcaptionnumbering{continuous}
\renewcommand{\lakitcaptionpunctuation}{:}

\definecolor{cBlue}{HTML}{000000}   \definecolor{cBlueBg}{HTML}{F2F2F2}
\definecolor{cGreen}{HTML}{000000}  \definecolor{cGreenBg}{HTML}{F2F2F2}
\definecolor{cAmber}{HTML}{000000}  \definecolor{cAmberBg}{HTML}{F2F2F2}
\definecolor{cPurple}{HTML}{000000} \definecolor{cPurpleBg}{HTML}{F2F2F2}
\definecolor{cTeal}{HTML}{000000}   \definecolor{cTealBg}{HTML}{FFFFFF}
\definecolor{cRed}{HTML}{000000}    \definecolor{cRedBg}{HTML}{F2F2F2}
\definecolor{cSlate}{HTML}{000000}  \definecolor{cSlateBg}{HTML}{F2F2F2}
\definecolor{cRule}{HTML}{000000}
\definecolor{cTblLine}{HTML}{767676}\definecolor{cTblBg}{HTML}{FFFFFF}

% The proven default mdframed renderer preserves both native directions.
% Its Q/A frames deliberately use square corners: its renderer does not draw
% roundcorner, while the TikZ alternative displaced complete frames in RTL.

% Mutable semantic-block hooks; provide+renew permits either preamble load
% order and also makes the theme useful without the semantic extension.
% The default renderer's leftline is the local paragraph's reading edge:
% physical right in Hebrew and physical left in English. Keep that native
% behavior inside language blocks as well as in either document main language.
\providecommand{\referencenotestyle}{}
\providecommand{\referencewarningstyle}{}
\providecommand{\referenceformulastyle}{}
\providecommand{\referenceqastyle}{}
\renewcommand{\referencenotestyle}{%
  linecolor=black,backgroundcolor=cAmberBg,
  leftline=true,rightline=false,topline=false,bottomline=false,
  roundcorner=0pt,innerleftmargin=3pt,innerrightmargin=3pt,
  innertopmargin=2pt,innerbottommargin=2pt,skipabove=3pt,skipbelow=3pt,
  startinnercode=\raggedright}
\renewcommand{\referencewarningstyle}{%
  linecolor=black,backgroundcolor=cRedBg,
  leftline=true,rightline=false,topline=false,bottomline=false,
  roundcorner=0pt,innerleftmargin=3pt,innerrightmargin=3pt,
  innertopmargin=2pt,innerbottommargin=2pt,skipabove=3pt,skipbelow=3pt,
  startinnercode=\raggedright}
\renewcommand{\referenceformulastyle}{%
  linecolor=black!65,backgroundcolor=white,
  leftline=true,rightline=true,topline=true,bottomline=true,
  roundcorner=0pt,innerleftmargin=3pt,innerrightmargin=3pt,
  innertopmargin=2pt,innerbottommargin=2pt,skipabove=3pt,skipbelow=3pt,
  startinnercode=\raggedright}
\renewcommand{\referenceqastyle}{%
  linecolor=black!65,backgroundcolor=cPurpleBg,
  leftline=true,rightline=true,topline=true,bottomline=true,
  roundcorner=0pt,innerleftmargin=3pt,innerrightmargin=3pt,
  innertopmargin=2pt,innerbottommargin=2pt,skipabove=3pt,skipbelow=3pt,
  startinnercode=\raggedright}
% Natural word spacing is easier to scan in a narrow reference column. In
% LuaLaTeX, rightskip follows the actual paragraph direction: \raggedright
% aligns English at physical left and Hebrew at physical right. It never
% changes pardir/bodydir or mirrors a mathematical object. The environment
% groups retain ordinary document justification outside reference columns.
\AtBeginEnvironment{referencecolumns}{%
  \raggedright
  \def\columnseprulecolor{\color{black!40}}}
% Compact scientific objects keep their measured core caption/anchor behavior.
% Reduce only this grouped reference-figure presentation's list/paragraph gaps;
% ordinary figures, standalone diagrams and tables retain their existing skips.
\AtBeginEnvironment{referencefigure}{%
  \setlength{\medskipamount}{.25\baselineskip}%
  \setlength{\topsep}{.2\baselineskip}%
  \setlength{\partopsep}{0pt}}

% Heading sizes follow the actual selected body size, including a compact
% layout's local size. Native \@startsection still owns counters, stars,
% bookmarks, marks and keeping each heading with the following paragraph.
\makeatletter
% Heading text sizes may be fractional, but the native Computer Modern math
% families have discrete design sizes. Choose those faces deliberately rather
% than asking NFSS for unavailable 8.775/7.41-point faces or sub-5-point scripts.
\newcommand{\hllpdf@referencemathface}[1]{%
  \@tempdima=#1pt\relax\def\hllpdf@referencemathchosen{5}%
  \ifdim\@tempdima>5.5pt\def\hllpdf@referencemathchosen{6}\fi
  \ifdim\@tempdima>6.5pt\def\hllpdf@referencemathchosen{7}\fi
  \ifdim\@tempdima>7.5pt\def\hllpdf@referencemathchosen{8}\fi
  \ifdim\@tempdima>8.5pt\def\hllpdf@referencemathchosen{9}\fi
  \ifdim\@tempdima>9.5pt\def\hllpdf@referencemathchosen{10}\fi
  \ifdim\@tempdima>10.475pt\def\hllpdf@referencemathchosen{10.95}\fi
  \ifdim\@tempdima>11.475pt\def\hllpdf@referencemathchosen{12}\fi
  \ifdim\@tempdima>13.2pt\def\hllpdf@referencemathchosen{14.4}\fi
  \ifdim\@tempdima>15.84pt\def\hllpdf@referencemathchosen{17.28}\fi
  \ifdim\@tempdima>19.01pt\def\hllpdf@referencemathchosen{20.74}\fi
  \ifdim\@tempdima>22.81pt\def\hllpdf@referencemathchosen{24.88}\fi}
\newcommand{\hllpdf@referencedeclaremathsize}[2]{%
  \@tempdima=\dimexpr#1pt*#2/100\relax
  \edef\hllpdf@referenceheadingmathsize{\strip@pt\@tempdima}%
  \hllpdf@referencemathface{\hllpdf@referenceheadingmathsize}%
  \let\hllpdf@referenceheadingmathtext\hllpdf@referencemathchosen
  \@tempdima=\dimexpr\hllpdf@referenceheadingmathsize pt*7/10\relax
  \edef\hllpdf@referenceheadingmathscript{\strip@pt\@tempdima}%
  \hllpdf@referencemathface{\hllpdf@referenceheadingmathscript}%
  \let\hllpdf@referenceheadingmathscript\hllpdf@referencemathchosen
  \@tempdima=\dimexpr\hllpdf@referenceheadingmathsize pt/2\relax
  \edef\hllpdf@referenceheadingmathscriptscript{\strip@pt\@tempdima}%
  \hllpdf@referencemathface{\hllpdf@referenceheadingmathscriptscript}%
  \DeclareMathSizes{\hllpdf@referenceheadingmathsize}%
    {\hllpdf@referenceheadingmathtext}{\hllpdf@referenceheadingmathscript}%
    {\hllpdf@referencemathchosen}}
\newcommand{\hllpdf@referencedeclareheadingmath}[1]{%
  \hllpdf@referencedeclaremathsize{#1}{135}%
  \hllpdf@referencedeclaremathsize{#1}{114}}
\forcsvlist{\hllpdf@referencedeclareheadingmath}%
  {5,6,6.5,7,7.2,8,9,10,10.95,11,12}
\@ifundefined{ReferenceBodyFontSize}{}{%
  \hllpdf@referencedeclareheadingmath{\ReferenceBodyFontSize}}
% A minipage resets paragraph glue; restore this theme's natural reading edge
% after entry capture starts, without changing language or direction.
\providecommand{\hllpdf@referenceentrysetup}{}
\renewcommand{\hllpdf@referenceentrysetup}{\raggedright}
% The optional layout invokes this presentation hook after selecting its local
% body font, so every skip follows the actual compact/readable baseline. It
% affects only that grouped reference scope, never ordinary document displays.
\providecommand{\hllpdf@referencebodysetup}{}
\renewcommand{\hllpdf@referencebodysetup}{%
  \setlength{\abovedisplayskip}{.35\baselineskip plus .1\baselineskip minus .07\baselineskip}%
  \setlength{\belowdisplayskip}{.35\baselineskip plus .1\baselineskip minus .07\baselineskip}%
  \setlength{\abovedisplayshortskip}{.12\baselineskip plus .1\baselineskip}%
  \setlength{\belowdisplayshortskip}{.2\baselineskip plus .1\baselineskip minus .05\baselineskip}%
  \setlength{\jot}{.18\baselineskip}}
% Match bold heading text only inside the optional scalable reference scope.
% A standalone theme retains its existing Computer Modern sizing/version.
\newcommand{\hllpdf@referenceheadingmath}{%
  \@ifundefined{ifhllpdf@referencemathactive}{}{%
    \ifhllpdf@referencemathactive\boldmath\fi}}
\newcommand{\hllpdf@referencelocalhead}[1]{%
  \fontsize{\dimexpr\f@size pt*135/100\relax}%
    {\dimexpr\f@size pt*155/100\relax}\selectfont
  \hllpdf@referenceheadingmath
  #1\par\nobreak\vskip 1pt\hrule height .6pt\relax}
% A genuinely wide topic uses the same native ruled heading at its full local
% linewidth. No invented colored family or separate numbering is introduced.
\newcommand{\hllpdf@referencewidehead}[1]{%
  \hllpdf@referencelocalhead{#1}}
\newcommand{\hllpdf@referencesecondaryhead}{%
  \fontsize{\dimexpr\f@size pt*114/100\relax}%
    {\dimexpr\f@size pt*136/100\relax}\selectfont
  \hllpdf@referenceheadingmath}
\newcommand{\hllpdf@referencehead}[1]{%
  \@ifundefined{ifhllpdf@referencewideactive}{%
    \hllpdf@referencelocalhead{#1}%
  }{%
    \ifhllpdf@referencewideactive
      \hllpdf@referencewidehead{#1}%
    \else
      \hllpdf@referencelocalhead{#1}%
    \fi}}
\renewcommand{\section}{\@startsection{section}{1}{\z@}%
  {-1.3ex\@plus-.3ex\@minus-.2ex}{.8ex\@plus.2ex}%
  {\normalfont\rmfamily\bfseries\color{black}\hllpdf@referencehead}}
\renewcommand{\subsection}{\@startsection{subsection}{2}{\z@}%
  {-1.1ex\@plus-.2ex\@minus-.2ex}{.5ex\@plus.1ex}%
  {\normalfont\rmfamily\bfseries\color{black}\hllpdf@referencesecondaryhead}}
\renewcommand{\subsubsection}{\@startsection{subsubsection}{3}{\z@}%
  {-1ex\@plus-.2ex\@minus-.2ex}{.4ex\@plus.1ex}%
  {\normalfont\rmfamily\bfseries\color{black}}}
% Report's native \chapter still owns the page transition, counter and marks.
% Only its typeset heading is compacted; starred chapters retain native
% numbering semantics. Article has no chapter formatter to override.
\@ifclassloaded{report}{%
  \renewcommand{\@makechapterhead}[1]{%
    \vspace*{8pt}%
    {\parindent\z@\normalfont\rmfamily\bfseries\color{black}%
     \ifnum\c@secnumdepth>\m@ne
       \@chapapp\space\thechapter\par\nobreak\vskip 3pt
     \fi
     \hllpdf@referencehead{#1}\nobreak\vskip 8pt}}%
  \renewcommand{\@makeschapterhead}[1]{%
    \vspace*{8pt}%
    {\parindent\z@\normalfont\rmfamily\bfseries\color{black}%
     \hllpdf@referencehead{#1}\nobreak\vskip 8pt}}%
}{}
\makeatother

% Ordinary semantic callouts use a pale title strip and a white reading field.
% An untitled callout remains white; the optional title is never invented or
% inherited from a surrounding frame. This is separate from the inline-title
% reference notes and atomic Q/A blocks defined above.
\mdfdefinestyle{hllpdfreferencetheme}{%
  frametitle={},leftline=true,rightline=true,topline=true,bottomline=true,
  backgroundcolor=white,
  linewidth=.6pt,roundcorner=0pt,
  innerleftmargin=4pt,innerrightmargin=4pt,
  innertopmargin=3pt,innerbottommargin=3pt,
  skipabove=5pt,skipbelow=5pt,everyline=true,
  startinnercode=\raggedright,
  splitbottomskip=3pt,splittopskip=5pt,
  frametitlebackgroundcolor=black!5,frametitlefontcolor=black,
  frametitlefont=\normalfont\rmfamily\bfseries,
  frametitleaboveskip=2pt,frametitlebelowskip=2pt}
\renewcommand{\hllpdfsemanticframestyle}{hllpdfreferencetheme}
\renewmdenv[style=hllpdfreferencetheme,linecolor=black]{defbox}
\renewmdenv[style=hllpdfreferencetheme,linecolor=black]{thmbox}
\renewmdenv[style=hllpdfreferencetheme,linecolor=black]{notebox}
\renewmdenv[style=hllpdfreferencetheme,linecolor=black]{exbox}
\renewmdenv[style=hllpdfreferencetheme,linecolor=black]{warnbox}
\renewmdenv[style=hllpdfreferencetheme,linecolor=black]{keybox}
% Keep the core table container's tested atomic/anchor behavior and existing
% codebox number geometry; only the palette and ordinary callouts change.
